\section*{Capítulo \ref{ch:b3:total-synthesis} --- Síntesis total: estrategia y rutas clásicas}

\begin{solution}{exo:b3:total-synthesis:1}
Al 90\,\%: $0.9^{10} = 35$\,\% y $0.9^{20} = 12$\,\%. Al 80\,\%: $0.8^{10} = 11$\,\% y $0.8^{20} = 1.2$\,\%.
\end{solution}

\begin{solution}{exo:b3:total-synthesis:2}
$5 + 3 + 1 + 4 = 13$ etapas en total; LLS: $5 + 1 + 4 = 10$.
\end{solution}

\begin{solution}{exo:b3:total-synthesis:3}
4-(Dimetilamino)butan-2-ona, $\ce{CH3COCH2CH2N(CH3)2}$: el enol de la acetona se adiciona al ion
iminio $\ce{CH2=N+(CH3)2}$.
\end{solution}

\begin{solution}{exo:b3:total-synthesis:4}
Se desconectan los dos enlaces del anillo en posición alílica respecto del C=C: buta-1,3-dieno y but-3-en-2-ona
(metil vinil cetona), que reaccionan al calentarlos.
\end{solution}

\begin{solution}{exo:b3:total-synthesis:5}
Lineal: $0.9^{21} = 11$\,\%. Convergente: $0.9^{11} = 31$\,\%, casi el triple.
\end{solution}

\begin{solution}{exo:b3:total-synthesis:6}
Ruta A: $(7 + 1)/12 = 67$\,\%. Ruta B: $7/9 = 78$\,\%.
\end{solution}

\begin{solution}{exo:b3:total-synthesis:7}
Sin ellos: $0.9^8 = 43$\,\%. Con cuatro etapas adicionales al 95\,\%: $43\,\% \times 0.95^4 = 35$\,\%.
\end{solution}

\begin{solution}{exo:b3:total-synthesis:8}
Primario: un silil éter voluminoso (\textit{terc}-butildifenilsililo), eliminado con fluoruro.
Secundario: un éter 4-metoxibencílico, eliminado por oxidación con DDQ. El alcohol
terciario, impedido, puede dejarse a menudo libre, o protegerse como un silil éter pequeño
escindido con ácido suave. Cada condición deja intactos los demás grupos.
\end{solution}

\begin{solution}{exo:b3:total-synthesis:9}
Las \omterm{def:b3:total-synthesis:total}{síntesis totales} necesitaban decenas de etapas con \omterm{def:b3:total-synthesis:architecture}{rendimientos globales} muy inferiores al 1\,\%; un
compuesto con el sistema de anillos completo y la mayoría de los estereocentros, extraído de las
acículas del tejo europeo, una fuente renovable, se convierte en unas pocas etapas.
\end{solution}

\begin{solution}{exo:b3:total-synthesis:10}
Adición de Michael: el enolato de la diona (C2, entre los carbonilos) se adiciona al
carbono $\beta$ de la metil vinil cetona, lo que da una tricetona. Aldólica: el enolato de la metil cetona
ataca intramolecularmente un carbonilo del anillo, y la deshidratación da la
ciclohexenona: la cetona de Wieland--Miescher, una enediona bicíclica con un grupo metilo
angular.
\end{solution}

\begin{solution}{exo:b3:total-synthesis:11}
$y^{m+1}/y^{2m+1} = y^{-m}$: al 90\,\%, 1.7 para $m = 5$ y 4.9 para $m = 15$. Cuanto más larga es la ruta, más
compensa la convergencia.
\end{solution}

\begin{solution}{exo:b3:total-synthesis:12}
Cada anillo se cierra con el catión y el doble enlace atacante en una disposición de tipo
silla; la adición a cada doble enlace es anti, y los sustituyentes adoptan
posiciones pseudoecuatoriales, lo que hace \textit{trans} todas las fusiones de anillo. Un doble enlace \textit{E}
sitúa sus dos continuaciones de cadena de modo que resulta la fusión \textit{trans}; uno \textit{Z}
daría la fusión \textit{cis}: la geometría del polieno codifica la estereoquímica del
producto.
\end{solution}

\begin{solution}{pb:b3:total-synthesis:1}
\textbf{1.} $\ce{C4H6O2}$, $\ce{CH5N}$, $\ce{C5H6O5}$, $\ce{C8H13NO}$.

\textbf{2.} $\ce{C4H6O2 + CH5N + C5H6O5 -> C8H13NO + 2CO2 + 2H2O}$.

\textbf{3.} 86.0, 31.0 y \qty{146.0}{g/mol}; tropinona: \qty{139.0}{g/mol}.

\textbf{4.} $139.0/(86.0 + 31.0 + 146.0) = 0.53$: 53\,\%.

\textbf{5.} Sus grupos $\ce{CH2}$ centrales, situados cada uno entre dos grupos carbonilo, se enolizan fácilmente en
agua; la propia acetona es demasiado poco nucleófila, y los grupos carboxilo activantes
se van después.

\textbf{6.} Los iones iminio necesitan amina libre (demasiado ácido la protona), y la
formación del enol y del iminio necesitan algo de ácido: un pH próximo a la neutralidad los equilibra.

\textbf{7.} El $\ce{CH3NH2}$ se adiciona a uno de los grupos aldehído; la pérdida de agua da el ion iminio
$\ce{R-CH=N+H-CH3}$.

\textbf{8.} Un carbono enólico del ácido acetonadicarboxílico ataca el carbono del iminio y forma un
enlace C--C y una amina secundaria.

\textbf{9.} La amina secundaria se condensa con el segundo grupo aldehído de la misma
cadena, lo que da un ion iminio cíclico (de cinco miembros).

\textbf{10.} El enol del otro lado de la cetona lo ataca, cerrando el
anillo de seis miembros y el esqueleto bicíclico.

\textbf{11.} Cada uno es un $\beta$-cetoácido, que pierde $\ce{CO2}$ a través de un estado de transición cíclico de seis
miembros y da el enol de la cetona.

\textbf{12.} Dos enlaces C--C y dos enlaces C--N.

\textbf{13.} $0.75^{14} = 0.018$: 1.8\,\%.

\textbf{14.} $42/1.78 = 24$.

\textbf{15.} 100\,\%: una sola etapa, toda ella de construcción de enlaces.

\textbf{16.} Una sola operación, agua como disolvente, componentes baratos, ningún grupo protector y
casi ningún residuo aparte del dióxido de carbono y el agua.

\textbf{17.} La planta construye el esqueleto del tropano a partir de un ion iminio cíclico
derivado de un aminoácido y de una unidad derivada del acetato, mediante el mismo tipo de química
de Mannich.

\textbf{18.} Las cascadas (y las \omterm{def:b3:pericyclic:families}{cicloadiciones} formadoras de anillos, como la reacción de
Diels--Alder).

\textbf{19.} No: un \omterm{def:b3:point-groups:mirror}{plano de simetría} pasa por el N, su grupo metilo, C3 y el oxígeno
carbonílico.

\textbf{20.} Cada uno tiene cuatro grupos distintos, pero son imágenes especulares uno del otro respecto del
propio \omterm{def:b3:point-groups:mirror}{plano de simetría} de la molécula: un compuesto de tipo meso.

\textbf{21.} Diastereómeros, que difieren en la orientación del OH de C3 respecto del
puente de nitrógeno.

\textbf{22.} No: el \omterm{def:b3:point-groups:mirror}{plano de simetría} se conserva en ambos.

\textbf{23.} Las dos caras del grupo carbonilo no son equivalentes: una mira hacia el
puente de nitrógeno, y la otra, hacia el puente de dos carbonos, que impiden de forma distinta.

\textbf{24.} La ruta en un solo recipiente da unas 24 veces el \omterm{def:b3:total-synthesis:architecture}{rendimiento global} de la ruta lineal de 14
etapas.
\end{solution}
